\documentclass[12pt]{article}
\usepackage[utf8]{inputenc}
\usepackage[T1]{fontenc}
\usepackage{hyperref}
\title{Guevea de Humboldt Zapotec Language Resources for K-12 Curriculum Development}
\author{Compiled by Claude Sonnet 5.5 (claude-sonnet-5-5) for Andr\'as Kornai}
\date{}
\begin{document}
\maketitle

\section{Introduction}
Guevea de Humboldt Zapotec (ISO 639-3 zpg, Glottocode guev1238), also called Guevea Zapotec or Northern Isthmus Zapotec, is spoken in Oaxaca, Mexico. Its vernacular name is Didz-ri\~e'. The 2000 census counted about 4,700 speakers. Wikitongues lists it as EGIDS 5 (developing) with a Latin script. The searches found no school materials, grammar or dictionary.

\section{Primary Educational Resources}
None were found.

\section{Linguistic and Reference Materials}
\begin{itemize}
\item Wikipedia article: \url{https://en.wikipedia.org/wiki/Guevea_Zapotec}
\item Numeral data collected in 2019 by SIL linguists (Max Planck channumerals database): \url{https://lingweb.eva.mpg.de/channumerals/Zapotec-Guevea-de-Humboldt.htm}
\item Wikitongues language page: \url{https://wikitongues.org/languages/zpg/}
\end{itemize}

\section{Religious and Community Resources}
No scripture or community organisation was found for this language specifically.

\section{Digital and Audio Resources}
\begin{itemize}
\item Global Recordings Network language page, which lists no recordings in this language: \url{https://www.globalrecordings.net/en/language/zpg}
\item Global Recordings Network audio for the related Guienagati variety (evangelism and Bible teaching): \url{https://glorecnet-website.k.globalrecordings.net/en/language/2828}
\end{itemize}

\section{Curriculum Development Context}
Published material is very thin. The existence of a Latin-script writing system and SIL numeral work suggests a starting point for orthography-based materials, but a curriculum would need new development with the community.
\end{document}